\section{Legends of War and Peace}

\textbf{I Ching Hexagram 26}

Ta Chu: The Taming Power of the Great:

Ken --- Keeping Still, Mountain

---------------------------

Ch'ien --- the Creative, Heaven

In life, there come certain moments when there is lot of strain and tension, when things seem to be going out of hand, our every movement causes a negative movement in the Cosmos. In \emph{The Vocation of Man} (1800), \textbf{Johann Gottlieb Fichte} says that

\begin{quotex}
you could not remove a single grain of sand from its place without thereby changing something throughout all parts of the immeasurable whole,
\end{quotex}


the renowned butterfly effect. When one focuses on self-development one should in such situations remember the inner strength and independence, and not let our own ego be drawn in conflicts and ``litigations'' with the ego of others, e.g. trying to succumb and defeat the other in something that might seem as an occult war. This is a challenge to our inner independence -- that will be lost if we are stubborn in our pursuit of what is right, moral and orderly or if on the other hand we are drawn in the litigations.

As the affairs of men with their settled rules sometimes do not follow the vertical morality or the Way, it is better for us to leave worldly affairs to those who choose to deal with them, in their own way as their own conscience guides them or misleads them. We have a cosmic obligation to ourselves, development of the virtue and determination in acquiring our inner independence which we choose to connect instead to the things Eternal. In our own anxiety, when each dawn heart wakes itself to choking pain... one has to stand still... still on the shore and let the big wave pass us by. Neither is the power of the sea lesser when it breaks on the rocky shore, nor is the rock weaker when hit by the waves. Staying calm, we acquire the way of Heaven as the 6th line of this hexagram shows.

One's creative energy should be restrained beneath the mountain of Keeping Still. Ta Ch'u focuses on the power of silence and stillness, controlling of our lower urges, desires, watching over our ego, harnessing the emotions under the yoke of discipline. While the other is running around demonstrating their influence over the multitude because he is endowed with Power, at the same time the one who is weak has no need to prove anything. Controlled Power has the effect as its hidden, occult, its influence of the will with the steadiness of a mountain allows it to become the tool of retrospection over events, personal and general. Certain power is accumulated in this way of stillness as well as foresight. Difficult burdens, it is known, take away also the power of appropriate expression, therefore silence. There is no necessity to feel insecurity about this, as the way of the Righteous is forgiveness and the way of the beast is known to us all, and the latter has been futile and doomed since the beginning. One has to have faith and trust in that. The I Ching teaches that Controlled Power results in return to Innocence -- all that is learned from patience and non-action when the will is in the service of acceptance and the good of all this.

\begin{quotex}
With coarse rice to eat, with water to drink, and my bended arm for a pillow -- I have still joy in the midst of these things. Riches and honours acquired by unrighteousness are to me as a floating cloud.
\flright{\textsc{Confucius}, \emph{The Confucian Analects}}
\end{quotex}


\paragraph{Legend of the Empty town}


The way of a trickster, deceit, Strength and Purpose.

\emph{Romance of the Three Kingdoms}, considered one of the four great classical novels of Chinese literature, is a 14th-century historical novel attributed to Luo Guanzhong, set in the turbulent years towards the end of the Han dynasty and the Three Kingdoms period in Chinese history, starting in 169 AD and ending with the reunification of the land in 280. The novel adapts and widens several ideas of Master Sun Tzu, that are canons in military and strategical tactics curriculum. In chapter 38, we study the strategy of the Empty town: Taoist Zhuge Liang (181-234), also known as Sleeping dragon, one of great Chinese military wizards meets Liu Bei (161-223) a claimant to the throne of the Han dynasty whose sun is setting at that time. This is Liu Bei's third visit to the Liang's hermit cottage, which he has never left. Liu Bei makes a deal with the Taoist who teaches him about the balance of power and secures his accession to the throne. One of the most renowned stories about master Zhuge Liang we read in the chapter 95, which takes place around year 228, five years after Liu Bei's death.

Master Zhuge is left behind in the Western city (Si Cheng) with a small army of 5000, in the time of the siege, when the city was attacked by an army of 50,000 from the north state of Wei, led by the respected and much feared marshal Sima Yi. Sleeping Dragon sent half of his men to move the grains and fodder from the town to a place where the majority of his forces were resettled. So he was left in the town with just 2500 of his men. His officers were horrified when they saw the enemy's army. Sleeping Dragon climbed onto the city walls and observed raising clouds of dust, glistening of metal, as he listened to the drumming of the Wei army unstoppable pace. He ordered that all military flags should be taken down, all men should stay on their respective positions, and that all 4 doors of the city should be opened wide and at each door 20 men should be placed in peasant clothes, sweeping the road, washing the cobbler stone, chattering, carrying things, i.e. making an illusion of a relaxed everyday life. Master dressed himself in his Taoist clothes adorned with finest crane feathers, he took an instrument that resembles a lute and climbed onto the watchtower above main city gates; he sat, lighted the incense and played a most peaceful melody.

In the meantime, the army of Sima Yi reached the town and, seeing this unexpected sight that does not resemble a siege preparedness, they did not dare to come in. When general Sima Yi saw the Taoist playing surrounded by a cloud of incense he got frightened, as he thought to himself -- how well he knows the Sleeping Dragon, and he is renowned for never taking risks, therefore opening of the doors must mean that he wanted to lure Sima Yi inside the city and come down on him with far more superior force. Thinking it must be a trap, Sima Yi ordered the return of his army. Sleeping Dragon's men thought that master made some magic with his instrument and so they asked him what he did. Master Zhuge explained that Sima Yi knows him as a man who never takes risks so he would suspect a trap as a first thing, and for that even the illusion of a trap would be enough. From this we conclude that master Zhuge has not gambled with faith, he had no choice than to create an illusion with the purpose of avoiding the actual physical battle.

\paragraph{Legend about Discipline}


\begin{quotex}
If you correct your mind, the rest of your life will fall into place.
\flright{\textsc{Lao Tzu}}
\end{quotex}


Most what we know about the persona of Sun Tzu comes to us from Sima Qian (ca 45-85 B.C.) from his later \emph{Records of the Grand Historian}. In his narrative we find this interesting story about discipline. Master Sun became renowned for his \emph{Art of War} and therefore he was invited for an audience with the King of Wu (514-496 B.C.). The King kindly asked Master Sun to demonstrate his skills for the courtiers and if it can be done with the court ladies, as it will be more fun. Master Sun complied with the wish of the King, as any good general would.

180 of the most beautiful court ladies were summoned. Master Sun divided them into two opposing armies, ahead of each company he designated the lieutenants amongst whom were two of the Kings most cherished concubines. Master Sun established the rules, which were simple enough for any person with average level of care to perform: Look and face front, look left, look right, look and face behind. The women confirmed that they understood the rules, Master Sun took an axe in his hands and started giving these orders, but the women giggled, some laughed freely and others even made jest of Master Sun, who stayed calm.

He said, most seriously, that if his orders were unclear, it is he, who is their general who should be blamed for it, and so he explained the rules again -- in all detail and with great patience, that the even a person with even less than average level of care and attention should be able to perform. He gave orders again with strict tone -- the ladies once more burst into laughter. Upon which Master Sun replied, that after he has explained the rules in detail and orders have now been made clear -- the guilt is not on him but on their lieutenants -- and therefore they should be executed. He gave an immediate order for them to be beheaded.

Suddenly the whole court including the King went silent -- as the lieutenants were two of the King's favourites whom he loved greatly. Therefore he sent a pageboy to Master Sun with a message that suchlike will not be necessary as the King might lose his appetite if it came to giving heed to Master Sun orders. Sun replied, that the King -- as the absolute authority in the Kingdom has given him power and share in this authority ( to be understood as the Trust given -- that one with such a charge will take care for the welfare of the Land and people to his utmost abilities, vigilantly and faithfully), and he (Sun) as appointed general will not fail this Trust by succumbing to certain orders of his Majesty that are not in concordance with previously stated, as well as such an act would serve as good example for the others. The King had nothing to reply to such an answer but affirmed that the ladies should be executed and two new ladies named lieutenants. After this, Master Sun gave his orders again: Look front, face front. There were no more ladies giggling, it was a perfect compliant army that listened and obeyed. Look left -- they looked left, in silence; look right -- they acted as one, and nobody could see where is the beginning of the line, where is its middle, and where is the rear. Upon finishing the demonstration, Sun bowed to the King and offered him a well-trained army that will comply with any order, go through fire and water for his Majesty. And the King was pleased.

Of course, this is a romanticised and somewhat exaggerated version of the events and persons; Chapter Nine of the \emph{Art of War} by Master Sun speaks in favour of that:

\begin{quotex}
Soldiers must be treated in the first instance with humanity, but kept under control by means of iron discipline. This is a certain road to victory. (verse 43.)
\end{quotex}


Yen Tzu B.C. 493 said of Su-ma Jang-chu: ``His civil virtues endeared him to the people; his martial prowess kept his enemies in awe.'' Cf. Wu Tzu, ch. 4 init.:

\begin{quotex}
The ideal commander unites culture with a warlike temper; the profession of arms requires a combination of hardness and tenderness. If a general shows confidence in his men but always insists on his orders being obeyed, the gain will be mutual. (verse 45.)
\end{quotex}


Tu Mu says:

\begin{quotex}
A general ought in time of peace to show kindly confidence in his men and also make his authority respected, so that when they come to face the enemy, orders may be executed and discipline maintained, because they all trust and look up to him.
\end{quotex}


What Sun Tzu has said in v. 44, however, would lead one rather to expect something like this:

\begin{quotex}
If a general is always confident that his orders will be carried out, etc.
\end{quotex}


\paragraph{Art of Peace}


Born in the state of Lu, 551 BC, Confucius unlike Master Sun never wished to partake in the King's power but he rather travelled the lands preaching the ethical doctrine of peace, goodness, self-restraint, righteousness, cult of ancestors, proper manners, and the power of the moral example in the hope that he will save the State from powers of chaos, without the necessity of using trickery and deceit, neither craftiness of the strategists nor draconian methods and laws that ended in violence, employed by some other masters. Confucius witnessed the crumbling of his civilisation, the end of an era where his world sank into meaningless barbarism and violence.

Out of his admirable teaching, ``100 schools of philosophy'' arose, folded time distance all the way to us, so that we too can benefit from this teaching of the peaceful way. It takes much discipline and practice. The Art of Peace is much harder to master than the Art of War. This mastery is a choice of rather few rulers in history, but their greatness lives eternally for we much rather remember and learn of those who kept Pax than about those who chose the way of the tyrant. Lao Tzu is right there -- one must harmonize one's microcosm with the Tao. War and violence are in opposition to the Way of the Tao. To Love one's neighbour. Nothing is lost in Love.


\flrightit{Posted on 2018-07-31 by Mu Nianci}

